%============================================================
%  LEX v3.0 — Template Respons Reviewer (Revise & Resubmit)
%  Perintah Copilot: /response-reviewer [nomor komentar]
%  Salin file ini dan isi sesuai komentar reviewer Anda
%============================================================

\documentclass[12pt,a4paper]{article}

%--- Paket utama ---
\usepackage[bahasa]{babel}
\usepackage[utf8]{inputenc}
\usepackage[T1]{fontenc}
\usepackage{times}
\usepackage[margin=2.5cm]{geometry}
\usepackage{setspace}
\onehalfspacing

%--- Format & Warna ---
\usepackage{xcolor}
\usepackage{mdframed}
\usepackage{enumitem}
\usepackage{booktabs}
\usepackage{array}
\usepackage{hyperref}
\hypersetup{colorlinks=true, linkcolor=black, urlcolor=blue}

%--- Kotak komentar reviewer ---
\mdfdefinestyle{reviewerbox}{%
  backgroundcolor = gray!10,
  linecolor       = gray!60,
  linewidth       = 1pt,
  leftmargin      = 0pt,
  rightmargin     = 0pt,
  innerleftmargin = 8pt,
  innerrightmargin= 8pt,
  innertopmargin  = 6pt,
  innerbottommargin=6pt,
  skipabove       = 6pt,
  skipbelow       = 4pt,
}

%--- Kotak respons penulis ---
\mdfdefinestyle{authorbox}{%
  backgroundcolor = blue!5,
  linecolor       = blue!50,
  linewidth       = 1pt,
  leftmargin      = 0pt,
  rightmargin     = 0pt,
  innerleftmargin = 8pt,
  innerrightmargin= 8pt,
  innertopmargin  = 6pt,
  innerbottommargin=6pt,
  skipabove       = 4pt,
  skipbelow       = 8pt,
}

%============================================================
\begin{document}

%------------------------------------------------------------
%  HEADER SURAT RESPONS
%------------------------------------------------------------
\begin{center}
  {\Large\bfseries Respons terhadap Komentar Reviewer}\\[0.5em]
  {\large [Judul Naskah]}\\[0.3em]
  \textit{Manuscript ID: [ID Naskah]} \quad Tanggal: \today
\end{center}

\bigskip

Kepada Yth.\\
Editor [Nama Jurnal]

\medskip

Terima kasih atas kesempatan untuk merevisi naskah kami. Kami telah menelaah seluruh komentar
reviewer dan melakukan revisi secara menyeluruh. Dokumen ini menyajikan respons terhadap
setiap komentar secara berurutan. Perubahan dalam naskah ditandai dengan warna biru
(\textcolor{blue}{teks revisi}) agar mudah dilacak.

\bigskip
\hrule
\bigskip

%============================================================
%  REVIEWER 1
%============================================================
\section*{Reviewer 1}

%------------------------------------------------------------
%  Komentar 1.1
%  Perintah Copilot: /response-reviewer 1.1
%------------------------------------------------------------
\noindent\textbf{Komentar 1.1}
\begin{mdframed}[style=reviewerbox]
  \textit{[Salin teks komentar reviewer di sini persis seperti aslinya.]}
\end{mdframed}

\noindent\textbf{Respons Penulis:}
\begin{mdframed}[style=authorbox]
  % USER: Jalankan /response-reviewer 1.1 di Copilot Chat untuk menghasilkan respons ini.
  % Copilot akan mengikuti struktur: tindakan revisi → lokasi perubahan → argumen penjelas (jika ada penolakan).

  Kami menerima komentar ini. Revisi dilakukan pada \textbf{Bagian [X], halaman [Y], paragraf [Z]}:

  \medskip
  \textcolor{blue}{\textit{``[Teks revisi yang ditambahkan/diubah dalam naskah]''}}

  \medskip
  [Penjelasan singkat mengapa revisi ini menjawab komentar reviewer.]
\end{mdframed}

%------------------------------------------------------------
%  Komentar 1.2
%------------------------------------------------------------
\noindent\textbf{Komentar 1.2}
\begin{mdframed}[style=reviewerbox]
  \textit{[Salin teks komentar reviewer di sini.]}
\end{mdframed}

\noindent\textbf{Respons Penulis:}
\begin{mdframed}[style=authorbox]
  % USER: Jalankan /response-reviewer 1.2 di Copilot Chat untuk menghasilkan respons ini.
\end{mdframed}

%------------------------------------------------------------
%  Tambahkan blok Komentar/Respons di atas untuk setiap
%  komentar Reviewer 1 yang tersisa
%------------------------------------------------------------

\bigskip
\hrule
\bigskip

%============================================================
%  REVIEWER 2
%============================================================
\section*{Reviewer 2}

%------------------------------------------------------------
%  Komentar 2.1
%------------------------------------------------------------
\noindent\textbf{Komentar 2.1}
\begin{mdframed}[style=reviewerbox]
  \textit{[Salin teks komentar reviewer di sini.]}
\end{mdframed}

\noindent\textbf{Respons Penulis:}
\begin{mdframed}[style=authorbox]
  % USER: Jalankan /response-reviewer 2.1 di Copilot Chat untuk menghasilkan respons ini.
\end{mdframed}

%------------------------------------------------------------
%  Tambahkan blok Komentar/Respons di atas untuk setiap
%  komentar Reviewer 2 yang tersisa
%------------------------------------------------------------

\bigskip
\hrule
\bigskip

%============================================================
%  RINGKASAN PERUBAHAN (Change Summary Table)
%  Bantu editor melacak semua revisi dalam satu tabel
%============================================================
\section*{Ringkasan Perubahan Naskah}

\begin{table}[h]
\centering
\begin{tabular}{@{} c p{3.2cm} p{3.2cm} p{4cm} @{}}
\toprule
\textbf{No.} & \textbf{Komentar} & \textbf{Lokasi Revisi} & \textbf{Tindakan} \\
\midrule
1.1 & [Ringkasan komentar] & Bagian X, hal.\ Y   & [Diterima / Ditolak — alasan] \\
\addlinespace
1.2 & [Ringkasan komentar] & Bagian X, hal.\ Y   & [Diterima / Ditolak — alasan] \\
\addlinespace
2.1 & [Ringkasan komentar] & Bagian X, hal.\ Y   & [Diterima / Ditolak — alasan] \\
\bottomrule
\end{tabular}
\caption{Tabel Ringkasan Respons terhadap Komentar Reviewer}
\end{table}

\bigskip

\noindent Hormat kami,\\[0.8em]
[Nama Penulis Korespondensi]\\
[Institusi]\\
[Email]

\end{document}
